\section{From a kernel relation to a solution}
\label{sec:kernel}

We now apply Proposition~\ref{prop:planar} to prove
Theorem~\ref{thm:main}. The required surfaces must be embedded in a disk:
an arbitrary pairing of variable occurrences would not guarantee
planarity. Following the construction in
\cite[Section~5]{OpenAIKervaire2026}, we obtain the bands as level arcs
of the variable-circle maps on a punctured disk. This also ensures that the coefficient map is
constant along every band.

\subsection{A punctured disk for the relation}

Let $g\in G$ lie in the normal closure of the relators. By definition,
there is a finite expression in $G*F_n$
\begin{equation}\label{eq:normal-product}
 g=\prod_{\alpha=1}^{u}
       b_\alpha w_{i_\alpha}^{\epsilon_\alpha}b_\alpha^{-1},
 \qquad b_\alpha\in G*F_n,\qquad\epsilon_\alpha\in\{1,-1\}.
\end{equation}
If $u=0$, then $g=1$. Assume $u>0$. Use the cyclic representatives
of the relators chosen in Section~\ref{sec:planar}; conjugating a
relator only changes the corresponding conjugators in
\eqref{eq:normal-product}.

Choose a based CW presentation space $Y$ with $\pi_1(Y)=G$, and put
\[
 Y'=Y\vee S^1_1\vee\cdots\vee S^1_n.
\]
The added circles represent the variables, so $\pi_1(Y')=G*F_n$.
Choose based loops in $Y$ for the coefficients and for $g$. No
asphericity or finiteness property of $Y$ is needed.

Let $\Delta$ be an oriented polygonal disk and remove the interiors
of $u$ disjoint polygonal disks $D_\alpha$ in its interior. On the
resulting punctured disk
\[
 P=\Delta\setminus\bigcup_{\alpha=1}^{u}\operatorname{int}D_\alpha
\]
there is a map $f:P\to Y'$ with these boundary values: the outer
boundary reads $g$ in $Y$, and $\partial D_\alpha$, oriented as the
boundary of the missing disk, reads
$w_{i_\alpha}^{\epsilon_\alpha}$. Each variable letter is represented
by one monotone traversal of its circle. The intervening coefficient
intervals are mapped to the chosen loops in $Y$; retain an interval
even when its label is the identity.

Here is an explicit reason that these values extend over $P$. Cut $P$
along disjoint arcs from the outer boundary to the holes, and map those
arcs to the loops for $b_\alpha$. Place their outer endpoints in a
constant portion of the outer loop, with the holes encountered in reverse
order after the outer traversal. The boundary of the cut disk reads
\[
 g\bigl(b_u w_{i_u}^{\epsilon_u}b_u^{-1}\bigr)^{-1}
  \cdots
  \bigl(b_1 w_{i_1}^{\epsilon_1}b_1^{-1}\bigr)^{-1}.
\]
It is nullhomotopic by \eqref{eq:normal-product}. A nullhomotopy
extends the map across the cut disk. The two copies of each cut arc
have the same map with opposite traversal, so the extension descends
to $P$. Notice that the induced boundary orientation of $P$ on a hole
is opposite to the missing-disk orientation; this explains the inverses
in the displayed word.

\subsection{Disjoint bands from circle levels}

Let $f_j:P\to S^1_j=\mathbb R/\mathbb Z$ be the projection of $f$
that collapses all the other wedge summands, and let $v:P\to Y$ be
the coefficient projection. We construct bands on which $v$ is exactly
constant, without changing the boundary values of any $f_j$.

Choose a common sufficiently fine finite triangulation of $P$ and
piecewise affine circle maps $\widetilde f_j$ satisfying
\[
 \widetilde f_j|_{\partial P}=f_j|_{\partial P},
 \qquad
 \operatorname{dist}(\widetilde f_j(p),f_j(p))<\frac1{12}
 \quad(p\in P).
\]
These approximations can be constructed directly. The boundary maps are
already piecewise affine in the angular coordinate. Subdivide to respect
their breakpoints and so finely that the image of each triangle under
each $f_j$ lies in a short arc of the circle. Lift that arc to $\mathbb R$
and interpolate its vertex values. The lifts on a shared edge differ
by an integer, so the interpolants agree as circle maps. Uniform
continuity gives the stated error bound.

For each $j$, choose $t_j\in(1/3,2/3)$ that is not the image of any
vertex under $\widetilde f_j$. The level set
$\widetilde f_j^{-1}(t_j)$ is a finite disjoint union of polygonal
circles and properly embedded polygonal arcs. On each triangle it
consists of line segments, and avoidance of vertices makes the segments
join without branching. Since the approximation is fixed on the
boundary, there is exactly one endpoint for each occurrence of
$x_j^{\pm1}$ on a hole and no endpoint on the outer boundary or on a
coefficient interval.

These level sets are disjoint for different indices. Indeed, at a point
of $\widetilde f_j^{-1}(t_j)$ the original $f_j$ is at distance more
than $1/4$ from the circle basepoint. Thus $f$ lands in the $j$th
variable circle away from the wedge point, and all its other circle
projections vanish. The same remains true in a neighborhood of the
level set. In particular, $v$ is constant at the basepoint on that
neighborhood. This argument uses the original wedge-valued map $f$;
the separately approximated maps need not themselves combine to a map
into the wedge.

Discard the closed level curves and choose thin, mutually disjoint
rectangular neighborhoods of the proper arcs inside these neighborhoods.
They are the bands. Each meets the holes in small intervals inside the
corresponding variable traversals, and every occurrence has exactly one
such interval. Coorientation of a circle level makes the two endpoint
crossing signs of an arc opposite with respect to the boundary
orientation of $P$. Reversing both orientations to the missing-disk
orientations leaves the signs opposite. Hence each band pairs the same
variable with opposite traversal signs, as required in
Proposition~\ref{prop:planar}.

Use the smaller attachment intervals as the slots. The pieces of a
variable traversal outside its slot project constantly to $Y$, so
shrinking the slots does not change any coefficient label between them.
Moreover, $v$ is constant on every band. Reinsert the disks $D_\alpha$
as the signed word disks. Their union with the bands is a compact
surface embedded in the interior of $\Delta$, with every slot paired.
Every connected component is planar, and its boundary words are exactly
the loops supplied by $v$ on those boundary circles.

\subsection{Filling from the inside out}

It remains to account for the possibility of several components nested
inside one another. Each connected component has one exterior boundary
circle whose Jordan disk contains the component; its other boundary
circles enclose the inner complementary disks. The exterior Jordan
disks of distinct components are disjoint or nested. Choose a component
whose exterior disk is minimal under inclusion. Its inner complementary
disks contain no other component, and hence contain no missing word disk.
The coefficient map $v$ is already defined on each of them. Their
boundary words are therefore trivial in $G$.

Proposition~\ref{prop:planar} now implies that the exterior boundary
word of the chosen component is also trivial. Its exact boundary loop
extends to a map of its Jordan disk into $Y$. Replace $v$ inside that
disk by this extension. The two maps agree on the boundary, so they
paste continuously. No other component lies in the disk, and the
coefficient labels and band constants on all remaining components are
unchanged.

Repeating this finite procedure removes all components. The result is
a map $\Delta\to Y$ extending the original outer loop for $g$.
Consequently $g=1$ in $G$. This proves the injectivity in
Theorem~\ref{thm:main}; its quotient supplies the simultaneous solution.
\qed

\begin{remark}\label{rem:infinite-systems}
The finite statement also implies the version with arbitrary sets of
variables and equations. Assume the exponent rows, each of finite
support, are linearly independent over $\mathbb Q$. Any element of
$G$ killed in the quotient has a normal-closure expression involving
only finitely many relators and finitely many variables, including those
in its conjugators. The corresponding finite exponent matrix still has
independent rows, so Theorem~\ref{thm:main} makes that element trivial.
\end{remark}
